\paragraph{The complex multiplication value comes from one elliptic curve.}
The period calculation needs no height formula for higher dimensional
CM abelian varieties. An elliptic curve, a local orthogonal decomposition,
and the product formula suffice. We retain the prescribed host condition:
$E$ is CM, $E/E^+$ is unramified at finite places, and
$F=\mathbf Q(\sqrt{-d})\subset E$, where $-d$ is the prescribed
odd fundamental discriminant. Put $m=[E:F]$.

\begin{lemma}\label{geo-pure-cm-local-basis}
Let $L$ be a projective $\mathcal O_E$-lattice of rank $p+q$,
with a skew-Hermitian form $h$, linear in its first argument,
taking values in
$\mathfrak D_{E/\mathbf Q}^{-1}$. Suppose
$\operatorname{Tr}_{E/\mathbf Q}h$ is a perfect alternating
form on $L$. At the half of the complex embeddings extending
one fixed embedding $F\hookrightarrow\mathbf C$, suppose
$i\tau h$ has constant signature $(p,q)$.
For every rational prime $\ell$ there are orthogonal $E$-subspaces
$W_+,W_-$, positive and negative at all these embeddings, and
vectors $x_1,\ldots,x_p\in L\cap W_+$ and
$x_{p+1},\ldots,x_{p+q}\in L\cap W_-$ which form a basis of
$L\otimes\mathbf Z_{(\ell)}$ over
$\mathcal O_E\otimes\mathbf Z_{(\ell)}$.
The intersections $L\cap W_\pm$ are saturated projective lattices;
their direct sum equals $L$ after this localization.
Writing $G_\pm$ for their Gram blocks, every $\ell$-adic embedding
$j$ satisfies
\[
 \log|j\det G_+|=-pD_\ell(j),\qquad
 \log|j\det G_-|=-qD_\ell(j),\qquad
 D_\ell(j)=\log|j\mathfrak D_{E/\mathbf Q}|.
\]
The assertion includes $\ell=2$ and empty blocks.
\end{lemma}
\begin{proof}
Write $A=\mathcal O_E\otimes\mathbf Z_{(\ell)}$.
Unramifiedness over $E^+$ gives a real generator $\delta_0$ of
the absolute different over $A$. Trace duality and perfection
show that $\delta_0h$ is a unimodular skew-Hermitian form on $L_A$.
There is an anti-invariant unit $u\in A$.
At an inert dyadic place take $a$ with unit residue trace and
$u=a-\bar a=2a-\operatorname{Tr}(a)$; at a split place choose
$a$ with residues $(1,0)$, so $a-\bar a$ has residues $(1,-1)$.
At the other inert places choose
$a-\bar a$ nonzero in the quadratic residue field.
Choose $a$ simultaneously by the Chinese remainder theorem and
set $u=a-\bar a$, preserving exact anti-invariance.
Thus $u\delta_0h$ is a unimodular Hermitian form.

It has an orthogonal basis over the semilocal ring $A$.
Indeed a nondegenerate Hermitian space over a quadratic finite
field extension has a vector of nonzero norm, also in characteristic
two: otherwise the values on $v+aw$ would force the nondegenerate
trace pairing to vanish. At a split factor choose mutually dual
vectors in the two factors. Lift these choices simultaneously,
split off the resulting unit norm line, and induct.
Clearing denominators prime to $\ell$ makes this basis $e_i$ integral.

Choose a positive $p$-frame at every complex embedding in the
specified half. In each coset $e_i+\ell L$, the Minkowski lattice
has bounded covering radius. Approximate a sufficiently large
positive scalar multiple of the prescribed frames. After rescaling,
the errors tend to zero, so the resulting $x_1,\ldots,x_p$ are
positive at every embedding. Their Gram matrix is still invertible
over $A$, since it is unchanged modulo every prime above $\ell$.
Consequently its inverse gives an integral orthogonal projector
$L_A\to\sum A x_i$.
Take $W_+$ to be their $E$-span and $W_-=W_+^\perp$.
The complementary lattice is free over $A$; clear denominators
prime to $\ell$ in a basis to obtain the remaining $x_i$.
The signature gives negativity on $W_-$.
The two Gram blocks of $u\delta_0h$ are unimodular, proving the
valuation formulas. Intersections with $E$-subspaces are saturated
and hence projective over the Dedekind domain $\mathcal O_E$.
For rank zero the determinant is one.
\end{proof}

\begin{proposition}\label{geo-cm-integral-product}
On $W_+$ use the fixed complex embedding of $F$, and on $W_-$
its conjugate. Give the resulting complex torus the Riemann
form $\operatorname{Tr}h$, and keep its specified symmetric
principal line bundle rigidified at zero.
It has a model over a number field, and after a finite extension
it has good reduction everywhere, with that same rigidified line.
For each chosen $\ell$, its tangent lattice is the exact image of
a product of integral elliptic tangent lattices under an isogeny
whose kernel is killed by an integer prime to $\ell$.
\end{proposition}
\begin{proof}
Choose a CM elliptic curve with period lattice homothetic to
$\mathcal O_F$. Use this one curve with the two conjugate
$\mathcal O_F$-actions. The pure vectors in the preceding lemma
give holomorphic homomorphisms from these curves to the torus.
Choose an $\mathcal O_F\otimes\mathbf Z_{(\ell)}$-basis of
$\mathcal O_E\otimes\mathbf Z_{(\ell)}$ and clear denominators
prime to $\ell$. Its products with the $x_i$ span a sublattice
of finite index prime to $\ell$. The product map $\psi$ is
therefore an isogeny, and some integer $N$ prime to $\ell$
kills its homology cokernel.

Polarized complex tori and their homomorphisms are algebraic
\cite[Theorems~2.8--2.9]{cm-milne-av}.
The finite kernel of $\psi$ is contained in the algebraic
$N$-torsion of the elliptic product, so the quotient descends to
$\overline{\mathbf Q}$
\cite[Remark~8.12]{cm-milne-av}.
Homomorphisms between abelian varieties over an algebraically
closed field do not increase after extending that field
\cite[Corollary~16.4]{cm-milne-av}; hence the coefficient action
and the principal polarization descend too. For
$\overline{\mathbf Q}\subset\mathbf C$, define the homomorphism
over a finitely generated regular subextension first.
To retain the specified line, first choose an algebraic line
inducing this polarization
\cite[Proposition~13.6, in characteristic zero]{cm-milne-av}.
Its Tate pairing is alternating because it becomes the prescribed
polarization after complex base change; one may choose an
auxiliary prime different from $2$.
Adjust it by a square root in $\operatorname{Pic}^0$ to make it
symmetric. Its difference from the specified symmetric complex
line is in $\operatorname{Pic}^0[2]$, hence is algebraic.
Choose the comparison to respect the zero rigidification.
All this data, and any prescribed finite level structure, descend
to one number field.

CM elliptic curves acquire good reduction everywhere after a
finite extension \cite[Theorem~12.14]{cm-milne-mf}.
Good reduction is invariant under isogeny
\cite[Theorem~1 and Corollary~2, p.~493]{cm-serre-tate},
so the quotient also has good reduction everywhere.
The Moret--Bailly theorem extends the specified symmetric
rigidified line after another finite extension
\cite[Theorem~3.4(i)]{pazuki-theta-faltings}.
Its identity component is the same abelian scheme, by the
N\'eron comparison already used in
Theorem~\ref{geo-classical-tangent-slope}.
The vertical divisor and zero rigidification argument there
preserves the line and its canonical metric exactly.

Finally $N\psi^{-1}$ sends period lattices into period lattices
and is an algebraic homomorphism $\chi$.
Both maps extend to the good models by the N\'eron mapping
property \cite[Definition~5.1 and Example~5.3]{cm-conrad-minimal}.
Their identities $\chi\psi=[N]$ and $\psi\chi=[N]$ extend.
The coefficient action and its polarization adjunction identities
extend by the same uniqueness. On products of the good model,
the zero rigidification again removes the vertical divisor in
the corresponding cubical line identifications.
At a place above $\ell$, differentiation gives inverse integral
maps $d\psi$ and $N^{-1}d\chi$. Thus the product tangent basis
is a basis of the prescribed integral lattice, without an index.
\end{proof}

\begin{lemma}\label{geo-cm-elliptic-period}
For a CM elliptic curve by $\mathcal O_F$, choose after finite
extension an everywhere good model and an integral tangent basis
$v$. In the metric $\|v\|_\sigma^2=\pi H_\sigma(v,v)$, put
$n_\sigma=\|v\|_\sigma^2$. Then
\[
 -\frac1{[M:\mathbf Q]}\sum_{\sigma:M\hookrightarrow\mathbf C}
       \log n_\sigma
 =x_d:=\gamma+\log\frac{2\pi}{\sqrt d}
                  -\frac{L'(1,\chi_{-d})}{L(1,\chi_{-d})}.
\]
Finite extension and reversal of the CM action preserve this mean.
\end{lemma}
\begin{proof}
The invariant differential line can be principalized after finite
extension: if an ideal $I$ satisfies $I^a=(b)$, adjoining an
$a$th root of $b$ makes its extension principal, since the group
of fractional ideals is torsion free. Let $\omega(v)=1$ and write
its period lattice as $\omega_1(\mathbf Z+\tau\mathbf Z)$,
with $y=\operatorname{Im}\tau>0$. The principal Riemann form gives
\[
 n_\sigma=\frac{\pi}{|\omega_1|^2y},\qquad
 \Delta_\omega(\sigma)
   =(2\pi)^{12}\omega_1^{-12}\eta(\tau)^{24}.
\]
The second formula is the classical discriminant product
\cite[Theorem~4.21]{cm-milne-mf}.
Good reduction and the integral differential make
$\Delta_\omega$ a unit at every finite prime, including $2$ and $3$
\cite[Example~5.3]{cm-conrad-minimal}.
Consequently
\[
 \log n_\sigma
   =-\log(4\pi y|\eta(\tau)|^4)
                         +\tfrac16\log|\sigma\Delta_\omega|,
\]
and the last term has mean zero.

The CM class theorem makes the ideal classes occur with equal
multiplicity among the embeddings
\cite[Theorem~12.18]{cm-milne-mf}.
For
$Z_\tau(s)=\sum' y^s/|m+n\tau|^{2s}$, the Kronecker limit formula is
\[
 Z_\tau(s)=\frac{\pi}{s-1}
       +2\pi\left(\gamma-\log2
                    -\log(\sqrt y|\eta(\tau)|^2)\right)+O(s-1)
 \quad\cite[\S2]{cm-zagier}.
\]
If $h_F$ is the class number and $w_F$ the number of roots of unity,
the ideal norm decomposition gives
\[
 \zeta_F(s)=\frac1{w_F}\left(\frac2{\sqrt d}\right)^s
                       \sum_{[\mathfrak a]}Z_{\tau_{\mathfrak a}}(s).
\]
Compare the constant term divided by the residue with
$\zeta_F(s)=\zeta(s)L(s,\chi_{-d})$. This gives
\[
 \frac1{h_F}\sum_{[\mathfrak a]}
       \log(4\pi y_{\mathfrak a}|\eta(\tau_{\mathfrak a})|^4)=x_d.
\]
Conjugation permutes the ideal classes and preserves this real
modular expression. The product formula and uniform class
multiplicities prove the asserted mean.
\end{proof}

\begin{proposition}\label{geo-cm-flat-degree}
Suppose the determinant family of
Lemma~\ref{geo-flat-determinant-bridge} is algebraically defined,
is constant in the period trivialization, and has the following
local norm compatibility: at each place its prescribed flat norm
at the point under comparison equals the norm at every allowed
CM point above, computed on the same integral eigenlattices and
canonical metric. This equality is independent of which CM
point is chosen at that place.
For constant signature $(p,q)$ its actual normalized degree is
\[
 \frac{p-q}{2}P_E,\qquad
 P_E=\gamma+\log(2\pi)-\log\operatorname{rd}(E)
                     -\frac{L'(1,\chi_{-d})}{L(1,\chi_{-d})}.
\]
If the finite logarithmic norms are corrected by subtracting
$q(\log|\mathfrak D_E|-\log|d|)$, the degree is instead
\[
             \frac{p-q}{2}P_E
                      +q(\log d-\log\operatorname{rd}(E)).
\]
These are the precise inputs of the calculation.
Theorem~\ref{geo-flat-family-generator} constructs the generic family
from the prescribed model, and Corollary~\ref{geo-flat-norm-assembly}
supplies the norm comparison from its exact integral identifications.
\end{proposition}
\begin{proof}
Fix an $E$-basis $e$ of the homology lattice, with Gram matrix $G_e$.
At a pure CM point write $x=Te$ and $t=\det T$.
For the algebraic eigenframes $v,w$ supplied by the elliptic product,
let $\varphi(v,w)$ denote the scalar of the determinant section.
The CM elliptic tangent vectors have period scalars $z_+,z_-$.
In the pure basis, the period map on the positive space and the
dual negative space has blocks
$z_+I_p$ and $z_-^{-1}G_-^{-1}$.
If $\rho$ is the constant in the flat period trivialization, then
\[
 \varphi(v,w)\tau(t)z_+^p=\rho z_-^q\tau(\det G_-).
\]
Thus
$r=\varphi(v,w)\tau(t)/\tau(\det G_-)$ is a nonzero algebraic
number and satisfies $rz_+^p=\rho z_-^q$.
This same $r$ applies at every CM point. One curve with opposite
actions can be used for both signs, so even $z_+=z_-$ is possible.

There is no extra metric constant in the product Gram matrix.
Indeed, with $\omega=(1+\sqrt{-d})/2$, the elliptic Hermitian
pairing is
\[
 \begin{aligned}
 h_F^\epsilon(x,y)
 &=\epsilon\sqrt{-d}\operatorname{Tr}_{E/F}h(x,y)\\
 &=\epsilon\bigl(B(\omega x,y)-B(x,y)+\omega B(x,y)\bigr),
 \qquad B=\operatorname{Tr}_{E/\mathbf Q}h.
 \end{aligned}
\]
Polarization adjunction therefore gives
$\langle d\psi_xv,d\psi_yv\rangle
=n_\sigma\,\sigma(h_F^\epsilon(x,y))$.
For an integral basis $b_i$ of $\mathcal O_E$ and its absolute
trace-dual $b_i^*$, the eigenprojection is
$P_\tau=\sum_i\tau(b_i^*)\rho(b_i)$.
The trace-dual identities collapse the two Gram sums, giving
\[
 \Gamma_+=n_+\epsilon_+\tau(\sqrt{-d}\,G_+),\qquad
 \Gamma_-=n_-\epsilon_-\bar\tau(\sqrt{-d}\,G_-).
\]
Here $\epsilon_\pm=\pm1$ and $\bar\tau=\tau\circ\overline{(\cdot)}$.
As
$\det G_+\det G_-=t\bar t\det G_e$, the complex logarithmic norm is
\[
 a_\sigma=\log|\sigma r|
 +\frac p2(\log n_{+,\sigma}+\tfrac12\log d)
 -\frac q2(\log n_{-,\sigma}+\tfrac12\log d)
 +\tfrac12\log|\sigma\tau(\det G_e)|.
\]

Fix a number field containing the given coefficient field
(in particular $\tau(E)$), the elliptic data and $r$.
At each $\ell$ choose the local pure basis in
Lemma~\ref{geo-pure-cm-local-basis}. Proposition
~\ref{geo-cm-integral-product} supplies its exact integral product
tangent basis. Proposition~\ref{geo-relative-different-eigenframe}
then supplies the integral eigenframes: multiply the positive
and negative projected frames by $\delta$ and $\bar\delta$,
where $\delta$ generates $\mathfrak D_{E/F}$ locally.
Using
$\log|j\delta|=\log|j\bar\delta|
=D_\ell(j)-\log|j\sqrt{-d}|$, their logarithmic norm is exactly
\[
 a_{\ell,j}
 =\log|jr|-\log|jt_\ell|-qD_\ell(j)
             +(q-p)(D_\ell(j)-\log|j\sqrt{-d}|).
\]
The transition determinant also satisfies
\[
 \log|jt_\ell|+\log|j\bar t_\ell|
 +\log|j\det G_e|=-(p+q)D_\ell(j).
\]
It has norm one at all embeddings for all but finitely many
$\ell$, because $e$ is then itself an integral local basis.

Average over all embeddings, with normalized absolute values.
Conjugation permutes the embeddings, so the two transition terms
have equal means. The product formula cancels $r$ and $\det G_e$.
The remaining three sums are
\[
 \overline{\log n_+}=\overline{\log n_-}=-x_d,\qquad
 \sum_\ell\overline{D_\ell}=-\log\operatorname{rd}(E),\qquad
 \sum_\ell\overline{\log|\sqrt{-d}|_\ell}=-\tfrac12\log d.
\]
The negative of the total logarithmic norm is consequently
$(p-q)(x_d+\tfrac12\log d-\log\operatorname{rd}(E))/2$.
This is the asserted degree. Subtracting the specified finite
correction increases the degree by
$q(\log d-\log\operatorname{rd}(E))$.
Each local CM calculation may use its own finite extension of
the fixed number field. Normalized means descend to that field;
there is no need for one extension containing every local CM point.
\end{proof}
